\chapter*{Wstęp}
\phantomsection
\addcontentsline{toc}{chapter}{Wstęp}
\label{chap:wstep}

Handel elektroniczny w ostatniej dekadzie stał się jednym z ~najszybciej rozwijających się
obszarów internetu. Systemy zamówień muszą mierzyć się z wieloma wyzwaniami.
Takimi jak zmienne w~czasie obciążenie, jak największa dostępność, krótkie czasy odpowiedzi oraz częstymi zmianami.
Architektura monolityczna w~której cała logika biznesowa zawiera się w~jednej aplikacji okazuje się niewystarczająca.
Zarówno pod względem skalowalności jak i~tempa rozwoju oprogramowania.

Odpowiedzią przemysłu informatycznego na te ograniczenia stała się architektura mikroserwisowa, 
polegająca na dekompozycji systemu na zbiór niewielkich, niezależnie wdrażanych usług, 
z~których każda odpowiada za wyodrębniony fragment domeny biznesowej i~komunikuje się z~pozostałymi 
poprzez dobrze zdefiniowane interfejsy \cite{newman2021}. Podejście to, choć rozwiązuje wiele problemów monolitu, 
wprowadza jednocześnie nowe wyzwania: konieczność zarządzania komunikacją między usługową, zapewnienia spójności 
danych rozproszonych pomiędzy usługami, a~także orkiestracji, monitorowania i~skalowania wielu niezależnych komponentów \cite{richardson2018}. 
W~praktyce przemysłowej wyzwania te adresowane są przy użyciu konteneryzacji, platform orkiestracji, z~spośród 
których standardem de facto stał się Kubernetes \cite{burns2022} oraz usług chmury publicznej, przejmujących 
od zespołów deweloperskich znaczną część odpowiedzialności za utrzymanie infrastruktury.

Niniejsza praca podejmuje się tego zagadnienia w~ujęciu praktycznym. Jej przedmiotem jest budowa kompletnego systemu 
obsługi zamówień zbudowanego. Utworzonego w~architekturze mikroserwisowej i~wdrożonego w~chmurze publicznej Amazon Web Services.
Z~wykorzystaniem zarządzanej usługi Kubernetes (Amazon EKS). Praca ta obejmuje pełny cykl budowy oprogramowania.
Zaczynając od analizy i projektu systemu. Idąc przez implementacje aplikacji webowej i jej backendu. Z wykorzystaniem
infrastruktury chmurowej w~podejściu \emph{infrastructure as code}. Kończąc na testach i analizie działania systemu
w~środowisku docelowym (chmurowym).

\section*{Cel i zakres pracy}
\phantomsection
\addcontentsline{toc}{section}{Cel i zakres pracy}
\label{sec:cel-i-zakres-pracy}

Celem tej pracy jest zaprojektowanie, implementacja oraz wdrożenie w~chmurze Amazon Web Services systemu obsługi
zamówień opartego na architekturze mikroserwisowej. Ponadto praca obejmuje weryfikację
poprawności działania systemu, która zostanie przeprowadzona za pomocą testów manualnych oraz obciążeniowych. Rezultatem tej pracy będzie wykazanie jak za pomocą wymienionych technologii
można wytworzyć relatywnie szybko system z dobrą skalowalnością, obserwowalnością oraz bezpieczeństwem.

Cel główny pracy można rozdzielić na:

\begin{itemize}
    \item teoretyczną analizę podstaw architektury mikroserwisowej oraz związanymi z nią technologiami (konteneryzacja, Kubernetes, usługi chmurowe AWS),
    \item zaprojektowanie modelu domenowego, architektury usług oraz kontraktów komunikacyjnych systemu obsługi zamówień, 
    \item implementację mikroserwisów realizujących obsługę zamówień, katalog produktów oraz proces płatności,
    \item implementację aplikacji webowej stanowiącej interfejs użytkownika systemu,
    \item budowę kompletnego środowiska uruchomieniowego w~chmurze AWS w~podejściu \emph{infrastructure as code},
    \item przeprowadzenie testów systemu oraz analizę jego działania i~kosztów utrzymania w~środowisku chmurowym.
\end{itemize}

Zakres pracy obejmuje wytworzenie następujących komponentów:

\begin{itemize}
    \item trzech mikroserwisów zaimplementowanych w~języku Go: serwisu zamówień (\code{order-service}), serwisu 
    produktów (\code{product-service}) oraz serwisu płatności (\code{payment-service}), komunikujących 
    się synchronicznie poprzez interfejsy REST/JSON oraz asynchronicznie poprzez zdarzenia domenowe publikowane do brokera Apache Kafka,
    \item aplikacji webowej typu SPA (\emph{single-page application}) zbudowanej w~oparciu o~framework Vue~3 i~język TypeScript, 
    umożliwiającej przeglądanie katalogu produktów, składanie zamówień oraz realizację płatności,
    \item warstwy persystencji opartej na relacyjnej bazie danych PostgreSQL, uruchamianej jako zarządzana usługa Amazon RDS,
    \item mechanizmów uwierzytelniania i~autoryzacji użytkowników zrealizowanych z~użyciem usługi Amazon Cognito i~tokenów JWT,
    \item infrastruktury chmurowej zdefiniowanej w~całości w~narzędziu Terraform,
    \item podsystemu monitorowania opartego na narzędziach Prometheus i~Grafana, uruchomionych wewnątrz klastra Kubernetes,
\end{itemize}

Poza zakresem pracy pozostaje implementacja usług pomocniczych jak powiadomienia oraz analityka. Integracja z~rzeczywistymi
operatorami płatności. Wytworzenie procesów ciągłej integracji za pomocą CI/CD zrealizowanych w GitHub Actions.
Tematy te zostaną omówione jedynie jako kierunki dalszego rozwoju systemu w~podsumowaniu pracy. 

\section*{Uzasadnienie wyboru tematu}
\phantomsection
\addcontentsline{toc}{section}{Uzasadnienie wyboru tematu}
\label{sec:uzasadnienie-wyboru-tematu}

Wybór tematu pracy wynika ze~znaczenia omawianych technologii we współczesnym przemyśle informatycznym.

Architektura mikroserwisowa wraz z~konteneryzacją i~orkiestracją przestały być rozwiązaniami niszowymi, 
zarezerwowanymi dla największych firm technologicznych. Badania społeczności Cloud Native Computing Foundation wskazują, 
że zdecydowana większość organizacji wykorzystujących technologie chmurowe stosuje Kubernetes w~środowiskach produkcyjnych, 
a~odsetek ten systematycznie rośnie \cite{cncf2023}. 

System zamówień jest bardzo dobrą dziedziną do wykorzystania architektury mikroserwisowej. Dzieli się ona łatwo na odrębne
tematy biznesowe takie jak: katalog produktów, zamówienia oraz płatności. Pomiędzy nimi występują interakcje synchroniczne
jak i asynchroniczne. Wymusza to zmierzenie się z~kluczowymi problemami systemów rozproszonych.

Literatura obszernie opisuje poszczególne technologie z~osobna, ale ich integracja w~spójny, działający system pozostaje
zadaniem wymagającym praktycznego doświadczenia. Zamiarem tej pracy było zweryfikowanie na rzeczywistym przykładzie, jakie kompromisy
i koszty niesie za sobą uruchomienie takiego systemu w~chmurze publicznej. Udokumentowani

\section*{Przegląd technologii chmurowych w kontekście mikroserwisów}
\phantomsection
\addcontentsline{toc}{section}{Przegląd technologii chmurowych w kontekście mikroserwisów}
\label{sec:przeglad-technologii-chmurowych}

Chmura obliczeniowa definiowana jest jako model umożliwiający wygodny, sieciowy dostęp na żądanie do współdzielonej puli konfigurowalnych 
zasobów obliczeniowych, które mogą być szybko przydzielane i~zwalniane przy minimalnym zaangażowaniu dostawcy \cite{nist2011}.
Można w niej wyróżnić trzy modele usługowe. SaaS (\emph{Software as a~Service}), w~którym produktem jest kompletne oprogramowanie dostępne przez sieć.
IaaS (\emph{Infrastructure as a~Service}), w~którym dostawca udostępnia surowe zasoby obliczeniowe, sieciowe i~dyskowe.
PaaS (\emph{Platform as a~Service}), w~którym udostępniana jest gotowa platforma uruchomieniowa dla aplikacji. Usługi wykorzystane w tej pracy
sytuują się pomiedzy modelami IaaS i~PaaS. Dostawca przyjmuje odpowiedzialność za utrzymanie platformy uruchomieniowej. Pozostawiając
użytkownikowi kontrolę nad jej konfiguracją. Z~perspektywy architektury mikroserwisowej kluczowe znaczenie mają cztery grupy technologii.

Konteneryzacja pozwalająca opakować usługę wraz z~jej zależnościami w~przenośny niezmienny obraz, której standardem de facto stał się Kubernetes: 
platforma automatyzująca rozmieszczanie, skalowanie, samonaprawianie i~aktualizacje kontenerów na klastrze maszyn \cite{burns2022}. 
Wszyscy najwięksi dostawcy chmury publicznej oferują zarządzane warianty Kubernetes.

\emph{infrastructure as code} (IaC), takie jak Terraform, pozwalające opisać całość infrastruktury chmurowej w~postaci deklaratywnego, wersjonowanego kodu.
Podejście to eliminuje ręczną, niepowtarzalną konfigurację środowisk i~umożliwia ich odtworzenie w~sposób w~pełni zautomatyzowany.
Ostatnią grupą są narzędzia wspierające eksploatację systemów rozproszonych. Takie jak systemy monitoringu i~obserwowalności.
Brokery komunikatów umożliwiające asynchroniczną wymianę zdarzeń. CI/CD automatyzujące zmiany w repozytorium i środowisku produkcyjnym.

\section*{Struktura pracy}
\phantomsection
\addcontentsline{toc}{section}{Struktura pracy}
\label{sec:struktura-pracy}

Praca składa się z~jedenastu rozdziałów, które można podzielić na część teoretyczną (rozdziały 1-4), projektową (rozdział 5), 
implementacyjno wdrożeniową (rozdziały 6-9) oraz badawczo-podsumowującą (rozdziały 10-11).

Rozdział~\ref{chap:podstawy-mikroserwisow} przedstawia teoretyczne podstawy architektury mikroserwisowej: porównanie z~architekturą monolityczną, 
modele komunikacji synchronicznej i~asynchronicznej oraz wzorce projektowe stosowane w~systemach rozproszonych. 
Rozdział~\ref{chap:kubernetes} omawia platformę Kubernetes: jej podstawowe obiekty, mechanizmy skalowania, trwałego składowania danych, 
obserwowalności i~bezpieczeństwa. Rozdział~\ref{chap:aws} charakteryzuje chmurę Amazon Web Services oraz usługi wykorzystane w~projekcie, 
wraz z~porównaniem Amazon EKS z~konkurencyjnymi zarządzanymi usługami Kubernetes.

Rozdział~\ref{chap:projekt-systemu} zawiera projekt systemu obsługi zamówień: opis funkcjonalny, model domenowy, architekturę usług, 
projekt interfejsów REST i~kontraktów zdarzeń, a~także uzasadnienie wyboru technologii implementacyjnych. 
Rozdziały~\ref{chap:implementacja-mikroserwisow} i~\ref{chap:implementacja-portalu} opisują implementację odpowiednio mikroserwisów 
backendowych oraz webowego portalu obsługi zamówień. Rozdział~\ref{chap:infrastruktura} przedstawia budowę kompletnego środowiska 
chmurowego w~podejściu \emph{infrastructure as code}.

Rozdział~\ref{chap:testy-i-analiza} poświęcony jest testom systemu oraz analizie jego działania w~środowisku docelowym, 
w~tym analizie kosztów wykorzystania usług chmurowych i~identyfikacji wąskich gardeł. Pracę zamyka rozdział~\enquote{\nameref{chap:wnioski}},
zawierający ocenę stopnia realizacji postawionych celów, omówienie napotkanych trudności oraz rekomendacje dotyczące dalszego rozwoju systemu.
